\section{Gleichförmige Bewegung}\label{sec:sample_motion}
Dieses Muster zeigt die Gestaltung von Definitionen, Beispielen und Abbildungen. Es ist keine vorgegebene Gliederung der Physikvorlesung.
\begin{DefinitionBox}{Geschwindigkeit}{sample_velocity}
Bei einer gleichförmigen Bewegung ist die \emph{Geschwindigkeit} \EnglishTerm{velocity} konstant. Für eine Bewegung entlang einer Achse gilt
\[
 v=\frac{\Delta s}{\Delta t},\qquad \Delta t>0.
\]
\end{DefinitionBox}
\begin{ExampleBox}{Gleichförmige Fahrt}{sample_motion}
Ein Wagen legt in \(\Delta t=\qty{4}{\second}\) eine Strecke von \(\Delta s=\qty{12}{\metre}\) zurück. Wir bestimmen seine Geschwindigkeit:
\begin{ClozeArea}[30]{3}
\[
 v=\frac{\Delta s}{\Delta t}=\frac{\qty{12}{\metre}}{\qty{4}{\second}}=\qty{3}{\metre\per\second}.
\]
\end{ClozeArea}
\begin{center}
\CachedTikzFigure{content/resources/chapters/sample/tikz/motion.tex}
\end{center}
\end{ExampleBox}
\begin{KeyPointBox}{Einheiten mitführen}{sample_units}
Die Einheit der Geschwindigkeit ergibt sich als Länge geteilt durch Zeit.
\end{KeyPointBox}
